\honest{Thou Art Physics}{Eliezer Yudkowsky}{2008}

Three months ago I set homework: trace the thinking that produces debates about free will, which is not the same as arguing whether free will exists. The wise reader guesses that it all adds up to normality, but how? Now people ask: ``If the future is determined, how can our choices control it?''

When people hear that physics is deterministic, their minds draw two causes, ``Me'' and ``Physics,'' each with an arrow to ``Future.'' If physics fixes the future, there is no room for me. This picture is an implicit representation in the brain that decides which arguments seem reasonable. \nb{No evidence is given about brains.} It is behind the neuroscience press releases saying ``it's not you who chooses, it's your brain.'' \nb{A Max Planck Society release two months earlier described the view that ``it is the brain that makes the decision, not a person's conscious mind.''} It is also behind ``that old chestnut'' that reductionism would make reasoning mere ``quarks bopping around.'' \nb{This is the argument from reason, best known from C.~S. Lewis, who called it a ``venerable philosophical chestnut.''}

The right diagram puts ``Me'' inside ``Physics.'' We miss it because no one can picture atoms making up a person the way one sees fingers making up a hand. There are many levels between our thoughts and quarks. Beliefs, desires, plans, emotions and temptations form the ``surface layer'' of the mind, visible without science: if I say that it is not you but your desires, plans and actions that determine the future, you see the part-whole relation at once, as with fingers and a hand. The relations below are not visible, so seeing yourself within physics takes ``constant vigilance,'' the point I said tripped up the quantum physicists. \nb{Von Neumann, for one, did apply quantum theory to measuring instruments.}

Compatibilism is the view that free will can be defined so as to fit deterministic physics; incompatibilism, that free will and determinism are incompatible. My position ``might perhaps be called'' Requiredism: agency, choice, control and responsibility ``require determinism,'' at least in patches, since you cannot plan and act ``amid utter chaos.'' \nb{That is a strong form of compatibilism, and an old one. Hume argued in 1748 that actions not caused by a person's character reflect neither credit nor blame on that person, and R.~E. Hobart's 1934 paper put the thesis in its title.}

Or perhaps the future must be ``determined by something.'' You do not need physics to push naive free will into incoherence: if the mind were not embodied in the brain, it would be embodied in something else, and if the future were not determined by physics, then by ``some law, some order'' that includes you. If physics controls us, how can we control ourselves? Turn it around: if physics did not control us, how could we? \nb{Agent-causal libertarians also hold that a free decision is determined by something, the agent, though not by prior events. The words ``some law, some order'' exclude them, and are not argued for.} How could thoughts judge other thoughts, emotions conflict, or uncertainty about plans become certainty ``in the midst of utter chaos?'' If we were not in reality, where could we be? \nb{Libertarian views in the literature do not posit chaos. Event-causal views such as Robert Kane's need only that some decisions be caused by the agent's reasons without being fixed by them. The post does not address such views.}

``The future is determined by physics,'' the kind of physics that includes human beings weighing decisions and being tempted. No quark ``swoops in from Pluto.'' Thoughts are real and made of smaller things, which we call physics. Physics includes our decisions. It does not explain them away. \nb{Elizabeth Anscombe gave C.~S. Lewis the same answer: a thought held for good reasons is rational, she wrote, ``whatever causal statements we make about him.''} Physics adds up to normality. The confusion comes from our cognitive algorithms.
